\section{Conclusion}

Predictive Memory Localization forecasts target, neighbor, capability, and
clean margin outcomes on a measured intervention grid. In the 3,000-record
primary study and 500-record confirmations across three base models, learned
directions improve target and clean incidence,
while the strength-disjoint low-dose response dominates static localization for
later-outcome prediction. A held-out selector improves utility over a fixed
coefficient, reduces neighbor damage, and replaces a dense scan with weak
probes plus action or abstention. PML therefore turns localization into a
falsifiable forecast and risk-aware decision: representation evidence identifies
leverage, while a small causal response reveals collateral risk before stronger
intervention.
